\chapter{Tests et validation}\label{chap:tests}

\section{Stratégie}
Deux niveaux se complètent.
\begin{enumerate}
  \item \textbf{Contrôles automatiques} : une classe \code{Verifications} lance \nbChecks{} vérifications Java, sans
        framework de test, depuis Eclipse (\textit{Run As, Java Application}). Le choix d'écrire de petites classes de contrôle
        plutôt que d'introduire JUnit garde le projet sans dépendance supplémentaire ; le chapitre~\ref{chap:perspectives}
        indique que JUnit serait la suite logique.
  \item \textbf{Parcours manuel} dans le navigateur, avec un scénario par règle de gestion (nominal et limite).
\end{enumerate}

\section{Contrôles automatiques}
Le tableau~\ref{tab:checks} indique ce que couvre chaque classe.

\begin{table}[htbp]
\centering
\caption{Classes de contrôle}\label{tab:checks}
\begin{tabularx}{\linewidth}{@{}lX@{}}
\toprule
\textbf{Classe} & \textbf{Contenu}\\
\midrule
\code{PanierCheck} & Règles du panier (ajout, quantités 1 à 99, produit inactif, stock, fusion de lignes, prix figé, total) et cohérence des relations entre entités.\\
\code{UtilChecks} & Hachage des mots de passe, formatage des montants, limitation des tentatives de connexion (avec horloge simulée).\\
\code{ServiceChecks} & Validation des formulaires (e-mail, mot de passe, référence, prix\ldots) et pagination.\\
\code{WebChecks} & Redirection de retour sûre (contre la redirection ouverte), jeton CSRF, JSON du mini-panier (échappement).\\
\code{JpaChecks} & Traduction des erreurs MySQL en messages métier et échappement des jokers du \code{LIKE}.\\
\bottomrule
\end{tabularx}
\end{table}

Ces contrôles ne se connectent pas à la base : ils testent la logique pure (entités, validation, hachage, JSON), ce qui les rend
rapides et reproductibles sur n'importe quel poste. Ce qui touche réellement MySQL (mapping, cascades, contraintes, verrou optimiste) est
vérifié autrement : par l'inspection directe de la base (chapitre~\ref{chap:persistance}) et par le parcours manuel ci-dessous.
Chaque contrôle affiche une ligne, et le bilan indique le nombre de succès et d'échecs ; un échec termine le programme avec un code
d'erreur.

\fig{18-verifications-console}{Résultat de la classe Verifications}{fig:verifications}

\section{Parcours manuel}
Le tableau~\ref{tab:manuel} liste les scénarios exécutés dans le navigateur avec les comptes de démonstration. Les comptes
sont \code{yasmine@example.com} et \code{omar@example.com} (clients, mot de passe \code{Client\#2026}) et \code{admin@example.com}
(mot de passe \code{Admin\#2026}).

\begin{table}[htbp]
\centering
\caption{Scénarios de test manuels}\label{tab:manuel}
\small
\begin{tabularx}{\linewidth}{@{}lXXl@{}}
\toprule
\textbf{N°} & \textbf{Action} & \textbf{Résultat attendu} & \textbf{Règle}\\
\midrule
T1 & Ouvrir la vitrine, choisir une catégorie, rechercher « clavier » & Liste filtrée, compteur à jour & ---\\
T2 & Omar ajoute un produit en stock & Tiroir ouvert, badge à 1, total correct & R4, R5\\
T3 & Omar ajoute le même produit une seconde fois & Une seule ligne, quantité 2 & R4\\
T4 & Omar ajoute \code{MAI-002} (stock 0) & Bouton désactivé ; une requête forcée est refusée & R3\\
T5 & Omar ajoute 4 unités de \code{INF-003} (stock 3) & Refus avec message de stock & R3\\
T6 & Omar saisit la quantité 0 puis 100 & Refus dans les deux cas & R1\\
T7 & Ajouter \code{MOD-003} (retiré de la vente) par requête forcée & Refus « plus disponible » & R2\\
T8 & Un visiteur non connecté ajoute au panier & Redirection vers la connexion & accès\\
T9 & Inscription avec une adresse déjà utilisée & Message d'erreur sous le champ & R7\\
T10 & Inscription avec un mot de passe « abcdefgh » & Refus (pas de chiffre) & validation\\
T11 & Connexion avec 6 mots de passe faux & Blocage temporaire au 6\textsuperscript{e} essai & force brute\\
T12 & Un client ouvre \code{/admin/produits} & Page 403 & accès\\
T13 & L'administrateur modifie le prix d'un produit présent dans un panier & Le prix de la ligne du panier ne change pas & R5\\
T14 & Supprimer un produit présent dans un panier & Refus, suggestion de le retirer de la vente & R9\\
T15 & Supprimer une catégorie qui contient des produits & Refus & R10\\
T16 & Deux onglets modifient le même produit ; enregistrer les deux & Le second est refusé (conflit de version) & verrou optimiste\\
T17 & Rafraîchir la page après un ajout & Aucun ajout en double (redirection après POST) & PRG\\
T18 & Envoyer un POST sans jeton (outil externe) & Refus 403 & CSRF\\
\bottomrule
\end{tabularx}
\end{table}

\fig{19-donnees-mysql-apres-action}{Contrôle dans MySQL après une action : la ligne de panier correspond à ce qui est affiché}{fig:mysqlapres}

\section{Ce que les tests ont révélé}
La première exécution sur Tomcat a montré que EclipseLink, à la demande de \code{create-or-extend-tables}, tentait de recréer des
contraintes existantes (codes MySQL 1061 et 1826). L'avertissement ne bloquait rien mais montrait un défaut de conception :
l'application ne doit pas toucher au schéma. Le réglage \code{jpa.ddl=none} est donc maintenant celui du fonctionnement normal.
